\section*{Lectures on Lie Groups and Lie Algebras – 1}

\textbf{Definition 1.}
A \emph{Lie algebra} is a vector space equipped with a bilinear operation 
$[\cdot,\cdot]$ (usually called the commutator) satisfying the following axioms:

\begin{enumerate}
\item (Skew-symmetry or anticommutativity) $[x,x] = 0$.
If the characteristic of the base field is not equal to $2$, this is equivalent to
\[
[x,y] = -[y,x].
\]

\item (Jacobi identity)
\[
[x,[y,z]] + [y,[z,x]] + [z,[x,y]] = 0.
\]
\end{enumerate}

\textbf{Examples.}

\begin{enumerate}
\item Let $A$ be an associative algebra. Define a Lie algebra $A^{-}$ as the vector space $A$
with operation
\[
[a,b] := ab - ba.
\]

\item The space $\mathbb{R}^3$ with the vector cross product is a Lie algebra over $\mathbb{R}$.

\item Any vector space with the identically zero commutator.
Such a Lie algebra is called \emph{abelian}.
\end{enumerate}

The following example is more general.

\textbf{Definition 2.}
A \emph{derivation} of an algebra $A$ is a linear operator
\[
D: A \to A
\]
satisfying the Leibniz identity:
\[
D(a\cdot b) = D(a)\cdot b + a\cdot D(b).
\]

Clearly, all derivations of $A$ form a vector space.

\textbf{Proposition 1.}
If $D_1, D_2$ are derivations of $A$, then
\[
D_1D_2 - D_2D_1
\]
is also a derivation of $A$.

Thus all derivations of $A$ form a Lie algebra, denoted $\mathrm{Der}(A)$.

\textbf{Example.}
Describe $\mathrm{Der}(\mathbb{C}[z])$, the Lie algebra of derivations of the polynomial algebra.
Let $D \in \mathrm{Der}(\mathbb{C}[z])$. Then
\[
D(1) = D(1\cdot 1) = 1\cdot D(1) + D(1)\cdot 1 = 2D(1),
\]
hence $D(1) = 0$.

By induction using the Leibniz rule,
\[
D(z^n) = n z^{n-1} D(z).
\]
Thus
\[
D = D(z)\frac{\partial}{\partial z}.
\]

Therefore $\mathrm{Der}(\mathbb{C}[z])$ is the space of polynomial vector fields
on the complex line with the commutator of vector fields.
This Lie algebra is usually denoted $W_1$.

\section*{Glossary}

\textbf{Definition 3.}
A vector subspace $I$ of a Lie algebra $L$ is called an \emph{ideal} if
for all $x \in I$ and $y \in L$,
\[
[x,y] \in I.
\]

\textbf{Proposition 2.}
The commutator operation is well-defined on the quotient space $L/I$.
The resulting Lie algebra is called the \emph{quotient algebra}.

\textbf{Definition 4.}
The \emph{commutator subalgebra} of $L$ is the subspace
\[
[L,L] \subset L
\]
spanned by elements $[x,y]$ for all $x,y \in L$.

\textbf{Proposition 3.}
$[L,L]$ is an ideal of $L$. The quotient $L/[L,L]$ is abelian.

\textbf{Definition 5.}
The \emph{center} of $L$ is
\[
Z(L) = \{x \in L \mid [x,y] = 0 \text{ for all } y \in L \}.
\]
Clearly $Z(L)$ is an ideal.

\textbf{Definition 6.}
Let $L_1, L_2$ be Lie algebras. A linear map
\[
\varphi: L_1 \to L_2
\]
is a \emph{homomorphism} if
\[
\varphi([x,y]) = [\varphi(x),\varphi(y)].
\]

The \emph{kernel} is
\[
\ker\varphi = \{x \in L_1 \mid \varphi(x)=0\},
\]
and the \emph{image} is
\[
\mathrm{Im}\,\varphi = \{\varphi(x) \mid x \in L_1\}.
\]

A bijective homomorphism is called an \emph{isomorphism}.

\textbf{Proposition 4 (Homomorphism Theorem).}
If $\varphi: L_1 \to L_2$ is a homomorphism, then
$\ker\varphi$ is an ideal of $L_1$,
$\mathrm{Im}\,\varphi$ is a subalgebra of $L_2$, and
\[
L_1/\ker\varphi \cong \mathrm{Im}\,\varphi.
\]

\section*{Lie Algebras of Small Dimension}

We classify complex Lie algebras of dimension $\le 3$ up to isomorphism.

\subsection*{Dimension 1}

Every 1-dimensional Lie algebra is abelian.
Basis $\{e_1\}$ with $[e_1,e_1]=0$.

\subsection*{Dimension 2}

Either:
\begin{itemize}
\item Abelian (unique), or
\item Non-abelian (unique up to isomorphism) with basis $\{e_1,e_2\}$ and
\[
[e_1,e_2]=e_1.
\]
\end{itemize}

\subsection*{Dimension 3}

There is a unique abelian algebra.

Case 1: $\dim [L,L] = 1$.

If $e_1 \in Z(L)$:
\[
[e_2,e_3] = e_1
\]
(Heisenberg algebra).

Otherwise:
\[
[e_1,e_2] = e_1, \quad [e_1,e_3]=[e_2,e_3]=0.
\]

Case 2: $\dim [L,L]=2$.

If $[L,L]$ is abelian:
\[
[e_1,e_3]=e_1, \quad [e_2,e_3]=\lambda e_2
\]
or Jordan form case
\[
[e_1,e_3]=e_1, \quad [e_2,e_3]=e_2 + e_1.
\]

Case 3: $[L,L]=L$.

Then $L$ is unique and isomorphic to $\mathfrak{sl}_2(\mathbb{C})$.

There exists a basis $\{h,e,f\}$ such that
\[
[h,e]=2e, \quad [h,f]=-2f, \quad [e,f]=h.
\]

Thus
\[
L \cong \mathfrak{sl}_2(\mathbb{C}),
\]
the Lie algebra of $2\times2$ traceless matrices with
\[
e=e_{12}, \quad f=e_{21}, \quad h=e_{11}-e_{22}.
\]